% ============================================================
%  Lernzettel Template (my_dhbw Stil, analysis.tex)
%  Eine Spalte, mit TOC, accent-Theme.
%  Renderer: pdflatex (zweimal)
% ============================================================
\documentclass[11pt, a4paper]{article}

\usepackage[ngerman]{babel}
\usepackage{fontspec}
\setmainfont[
  Path=/usr/share/fonts/TTF/,
  Extension=.ttf,
  BoldFont=DejaVuSans-Bold,
  ItalicFont=DejaVuSans-Oblique,
  BoldItalicFont=DejaVuSans-BoldOblique
]{DejaVuSans}
\setsansfont[
  Path=/usr/share/fonts/TTF/,
  Extension=.ttf,
  BoldFont=DejaVuSans-Bold,
  ItalicFont=DejaVuSans-Oblique,
  BoldItalicFont=DejaVuSans-BoldOblique
]{DejaVuSans}
\setmonofont[Path=/usr/share/fonts/TTF/,Extension=.ttf]{DejaVuSansMono}
\usepackage[top=2cm, bottom=2cm, left=2cm, right=2cm, footskip=1cm]{geometry}
\usepackage{amsmath, amssymb}
\usepackage{xcolor}
\usepackage{titlesec}
\usepackage{titletoc}
\usepackage{enumitem}
\usepackage{fancyhdr}
\usepackage{lastpage}
\usepackage{microtype}
\usepackage{parskip}

\usepackage{booktabs}
\usepackage{array}
\usepackage{longtable}
\usepackage{tabularx}
\usepackage{ltablex}
\keepXColumns
\usepackage{colortbl}
\usepackage[most,breakable]{tcolorbox}
\usepackage{listings}
\usepackage[normalem]{ulem}
\usepackage{pdflscape}
\usepackage[colorlinks=true, linkcolor=accent, urlcolor=accent]{hyperref}

\renewcommand{\familydefault}{\sfdefault}

% ── Theme (my_dhbw accent) ────────────────────────────────────
\definecolor{accent}{RGB}{0, 60, 113}
\definecolor{primaryblue}{RGB}{0, 60, 113}
\definecolor{codebg}{RGB}{245, 245, 245}
\definecolor{figurebg}{RGB}{232, 241, 255}
\definecolor{tablehintbg}{RGB}{240, 255, 240}
\definecolor{solutionbg}{RGB}{242, 255, 242}
\definecolor{rulegray}{RGB}{180, 180, 180}
\definecolor{tableheadbg}{RGB}{0, 60, 113}

% ── Header/Footer ─────────────────────────────────────────────
\pagestyle{fancy}
\fancyhf{}
\renewcommand{\headrulewidth}{0pt}
\renewcommand{\footrulewidth}{0.4pt}
\renewcommand{\sectionmark}[1]{\markboth{#1}{}}
\fancyfoot[L]{\footnotesize\color{accent}\nouppercase{\leftmark}}
\fancyfoot[R]{\footnotesize\color{accent}\thepage\,/\,\pageref{LastPage}}
\fancypagestyle{plain}{%
  \fancyhf{}%
  \renewcommand{\headrulewidth}{0pt}%
  \renewcommand{\footrulewidth}{0.4pt}%
  \fancyfoot[L]{\footnotesize\color{accent}\nouppercase{\leftmark}}%
  \fancyfoot[R]{\footnotesize\color{accent}\thepage\,/\,\pageref{LastPage}}%
}

% ── Typografie ────────────────────────────────────────────────
\titleformat{\section}
    {\color{accent}\large\bfseries}{}{0pt}{}
    [\vspace{-4pt}\color{accent}\rule{\linewidth}{2pt}]
\titleformat{\subsection}
    {\normalsize\bfseries\color{accent!85!black}}{}{0pt}{}
    [\vspace{-3pt}\color{accent!40}\rule{\linewidth}{0.4pt}]
\titleformat{\subsubsection}
    {\small\bfseries\color{accent!70!black}}{}{0pt}{}{}
\titlespacing*{\section}{0pt}{12pt}{4pt}
\titlespacing*{\subsection}{0pt}{8pt}{2pt}
\titlespacing*{\subsubsection}{0pt}{6pt}{1pt}

\setlength{\parindent}{0pt}
\setlength{\parskip}{4pt}
\renewcommand{\arraystretch}{1.25}
\setlist[itemize]{noitemsep, topsep=3pt, parsep=0pt, leftmargin=1.4em}
\setlist[enumerate]{noitemsep, topsep=3pt, parsep=0pt, leftmargin=1.6em}
\setlist[itemize,2]{label=\textendash, leftmargin=1.4em}

% ── Boxen ──────────────────────────────────────────────────────
\newtcolorbox{figurebox}[1]{
  colback=figurebg, colframe=accent!50,
  title={Abbildung: #1}, fonttitle=\small\bfseries,
  breakable, before skip=6pt, after skip=6pt
}
\newtcolorbox{tablehintbox}[1]{
  colback=tablehintbg, colframe=green!50!black!50,
  title={Tabelle: #1}, fonttitle=\small\bfseries,
  breakable, before skip=6pt, after skip=6pt
}
\newtcolorbox{solutionbox}[1]{
  colback=solutionbg, colframe=green!60!black!60,
  title={#1}, fonttitle=\small\bfseries,
  breakable, before skip=6pt, after skip=6pt
}

\lstset{
  basicstyle=\ttfamily\small,
  backgroundcolor=\color{codebg},
  breaklines=true,
  frame=single, framerule=0pt,
  xleftmargin=4pt, xrightmargin=4pt,
  aboveskip=6pt, belowskip=6pt,
  keepspaces=true
}

% ── TOC-Look ──────────────────────────────────────────────────
\renewcommand{\contentsname}{\color{accent}Inhalt}

\providecommand{\nichttitle}{Lernzettel}
\providecommand{\nichtsubtitle}{}

\renewcommand{\nichttitle}{Analysis 2}
\renewcommand{\nichtsubtitle}{Lernzettel}


\begin{document}

\begin{center}
{\bfseries\Large\color{accent}\nichttitle}\\[2pt]
\color{accent}\rule{0.35\linewidth}{1pt}\color{black}
\end{center}
\vspace{2pt}

{\small\tableofcontents}
\newpage

% ── BODY ──────────────────────────────────────────────────────

\section{Grundlagen Differenzialrechnung}


\subsection{Stetigkeit}
\textbf{Stetigkeit bei x₀}: lim(x→x₀) f(x) = f(x₀) — kein Sprung/keine Lücke.

\begin{itemize}
  \item Linksseitiger = rechtsseitiger Grenzwert = f(x₀)
  \item Beispiel 1/x: nicht stetig bei x₀=0
\end{itemize}


\subsection{Differenzierbarkeit}
\textbf{Differenzierbar bei x₀}: f'(x₀) = lim(h→0) [f(x₀+h)−f(x₀)]/h existiert.

\begin{itemize}
  \item Setzt Stetigkeit voraus; kein Sprung/Lücke/Knick
  \item Stetig-diffbar: zusätzlich keine Sprünge/Lücken in f'
\end{itemize}


\subsection{Ableitungsregeln}
\begin{itemize}
  \item \textbf{Potenz}: (a·xᵇ)' = b·a·xᵇ⁻¹
  \item \textbf{Logarithmus}: (a·ln x)' = a/x; (ln u)' = u'/u
  \item \textbf{Exponential}: (a·eˣ)' = a·eˣ
  \item \textbf{Summenregel}: termweise
  \item \textbf{Produktregel}: (gh)' = g'h + gh'
  \item \textbf{Quotientenregel}: (g/h)' = (g'h − gh')/h²
  \item \textbf{Kettenregel}: g(h(x))' = g'(h(x))·h'(x)
\end{itemize}


\section{Partielle Ableitungen}

\textbf{Partielle Ableitung ∂f/∂x}: Ableitung nach einer Variable, andere als Konstanten behandeln. Konstante Summanden ohne Differenzierungsvariable entfallen.

\begin{itemize}
  \item 1/z = z⁻¹; √(...) = (...)\textasciicircum{}0.5
\end{itemize}


\subsection{Beispiele}

{\small
\begin{tabularx}{\linewidth}{XXXX}
\hline
\rowcolor{tableheadbg}
{\color{white}\textbf{f}} & {\color{white}\textbf{∂f/∂x}} & {\color{white}\textbf{∂f/∂y}} & {\color{white}\textbf{∂f/∂z}} \\[2pt]
\hline
\endhead
\hline
\endfoot
2x³+y²+xy & 6x²+y & 2y+x & — \\
(x+y²)/z & 1/z & 2y/z & −(x+y²)/z² \\
x·ln(y)/(yz) & ln(y)/(yz) & (x/z)·(1−ln y)/y² & −x·ln(y)/(yz²) \\
2x·ln(xy²z) & 2ln(xy²z)+2 & 4x/y & 2x/z \\
\end{tabularx}
}



\section{Vektoranalysis}


\subsection{Skalarfeld vs. Vektorfeld}

{\small
\begin{tabularx}{\linewidth}{XXX}
\hline
\rowcolor{tableheadbg}
{\color{white}\textbf{Typ}} & {\color{white}\textbf{Abbildung}} & {\color{white}\textbf{Beispiel}} \\[2pt]
\hline
\endhead
\hline
\endfoot
Skalarfeld & ℝⁿ → ℝ & f(x,y)=10−x²+y² \\
Vektorfeld & ℝᵐ → ℝⁿ & (−y/(x²+y²), −x/(x²+y²))ᵀ \\
\end{tabularx}
}



\subsection{Vektornorm (euklidisch)}
\textbf{|x|} = √(x₁²+…+xₙ²)



\subsection{Totales Differenzial}
\textbf{df = (∂f/∂x)dx + (∂f/∂y)dy} — lineare Approximation der Funktionsänderung.



\subsection{Richtungsableitung}
Ableitung nach \textbf{normiertem} Vektor a mit |a|=1.

\begin{itemize}
  \item Diagonale dx=dy → dx=dy=√0.5
  \item dy=√3·dx → dx=0.5, dy=0.5√3, a=(0.5; 0.5√3)
\end{itemize}


\subsubsection{Beispiele df / Richtungswerte}

{\small
\begin{tabularx}{\linewidth}{XXXXXX}
\hline
\rowcolor{tableheadbg}
{\color{white}\textbf{f}} & {\color{white}\textbf{df}} & {\color{white}\textbf{Stelle}} & {\color{white}\textbf{Rtg. x}} & {\color{white}\textbf{Rtg. y}} & {\color{white}\textbf{Rtg. a}} \\[2pt]
\hline
\endhead
\hline
\endfoot
x²ln(y) & 2x·ln(y)dx + (x²/y)dy & (1,e) & 2 & 1/e & 1+√3/(2e) \\
x·eʸ & eʸdx + x·eʸdy & — & — & — & — \\
√(xy) & (√y/(2√x))dx + (√x/(2√y))dy & (4,0) & 0.25 & 2 & 0.125+√3 \\
√(xy) & dito & (4,4) & 0.5 & 0.5 & 0.25(1+√3) \\
x/y² & (1/y²)dx − (2x/y³)dy & (1,2) & 0.5 & −0.25 & 0.25−0.125√3 \\
x/y² & dito & (1,1) & 1 & −2 & 0.5−√3 \\
\end{tabularx}
}



\subsection{Gradient}
\textbf{grad(f) = ∇f = (∂f/∂x₁, …, ∂f/∂xₙ)ᵀ} — Vektor aller partiellen Ableitungen.

\begin{itemize}
  \item Zeigt in Richtung des steilsten Anstiegs
  \item Länge = Betrag des Anstiegs
  \item Notwendige Bedingung Extremstelle: grad(f)=0
\end{itemize}


{\small
\begin{tabularx}{\linewidth}{XXX}
\hline
\rowcolor{tableheadbg}
{\color{white}\textbf{f(x,y)}} & {\color{white}\textbf{grad(f)}} & {\color{white}\textbf{Stelle/Wert}} \\[2pt]
\hline
\endhead
\hline
\endfoot
x²+eʸ & (2x, eʸ) & (2, ln3) → \textbackslash{} \\
x²y−y³ & (2xy, x²−3y²) & (2,2) → ≈11.3 \\
xy+x+y & (y+1, x+1) & (5,7) → 10 \\
x²y−y & (2xy, x²−1) & grad=0 bei x=±1 ∧ y=0 \\
2−x−2y+xy & (−1+y, −2+x) & grad=0 bei (2,1) \\
xy²−y & (y², 2xy−1) & keine Nullstelle \\
x²/(2y) & (x/y, −x²/(2y²)) & grad=0 bei x=0 ∧ y≠0 \\
y+(y−1)x & (y−1, 1+...) & x=−0.5, y=1 \\
\end{tabularx}
}



\subsection{Divergenz}
\textbf{div(f) = ∇·f = ∂f₁/∂x + ∂f₂/∂y} — Skalarprodukt von ∇ mit Vektorfeld; misst Quellenstärke.

\begin{itemize}
  \item div>0: Quelle | div<0: Senke | div=0: quellenfrei
\end{itemize}


{\small
\begin{tabularx}{\linewidth}{XXX}
\hline
\rowcolor{tableheadbg}
{\color{white}\textbf{f}} & {\color{white}\textbf{div(f)}} & {\color{white}\textbf{quellenfrei wenn}} \\[2pt]
\hline
\endhead
\hline
\endfoot
(1,1) & 0 & ✓ überall \\
(−y, x) & 0 & ✓ überall \\
(−x, y²) & −1+2y & y=0.5 \\
(−y/(x²+y²), −x/(x²+y²)) & −1/(x²+y²) & nicht möglich \\
(x−y, x+y) & 2 & nicht möglich \\
(x²y², xy) & 2xy²+x & x=0 \\
(√x, y²) & 1/(2√x)+2y & y=−1/(2√x) \\
(yeˣ, ln(y)) & yeˣ+1/y & nicht möglich für x∈ℝ \\
(x²y, x−y²) & 2xy−2y & y=0 ∨ x=1 \\
\end{tabularx}
}



\subsection{Rotation}
\textbf{rot(f) = ∇×f} — Kreuzprodukt von ∇ mit Vektorfeld; misst Wirbelstärke.

\begin{itemize}
  \item rot=0: wirbelfrei
  \item ℝ²: Skalar ∂f₂/∂x − ∂f₁/∂y
  \item ℝ³: Vektor (∂g₃/∂y−∂g₂/∂z, ∂g₁/∂z−∂g₃/∂x, ∂g₂/∂x−∂g₁/∂y)
\end{itemize}


{\small
\begin{tabularx}{\linewidth}{XXX}
\hline
\rowcolor{tableheadbg}
{\color{white}\textbf{f}} & {\color{white}\textbf{rot(f)}} & {\color{white}\textbf{wirbelfrei wenn}} \\[2pt]
\hline
\endhead
\hline
\endfoot
(1,1) & 0 & ✓ \\
(−y, x) & 2 & nie \\
(−x, y²) & 0 & ✓ \\
(yeˣ, xeʸ) & eʸ−eˣ & x=y \\
(y, √x) & 1/(2√x)−1 & x=0.25 \\
(x²y², xy) & y−2x²y & x=±√0.5 ∨ y=0 \\
(x+y, y/x) & −y/x²−1 & y=−x² \\
(x²y, x−y²) & 1−x² (Skript) & x=±1 \\
(xyz, xy, x) (3D) & (−xy, xyz−1, y−xz) & — \\
\end{tabularx}
}



\subsection{Laplace-Operator}
\textbf{∆f = ∇·∇f = div(grad(f))} — Anwendung: Bild-/Kantenerkennung (diskret).



\subsection{Differenzialoperator-Identitäten}
\begin{itemize}
  \item \textbf{rot(grad(f)) = ∇×∇f = 0} (Satz von Schwarz)
  \item \textbf{div(rot(f)) = 0} für f(x,y,z) = (f₁(x,y), f₂(x,y), 0)ᵀ
\end{itemize}


\subsection{Relevanz Data Science}
\begin{itemize}
  \item \textbf{Lineare Regression}: yᵢ = β₀ + β₁xᵢ,₁ + … + ε
  \item \textbf{Loss (RMSE)}: f(β) = √((1/n)·Σ(yᵢ−ŷᵢ)²)
  \item \textbf{Gradient Descent}: Schritt entlang −grad(f); rot(grad)=0 → keine Endlosschleife
  \item \textbf{Diskretisierung}: Schrittweite, lokale vs. globale Minima nicht unterscheidbar
\end{itemize}


\section{Trigonometrie}


\subsection{Bogenmaß}
360°=2π, 180°=π, 90°=π/2, 45°=π/4 (Einheit "rad" oft weggelassen).



\subsection{Trig. Funktionen}
\begin{itemize}
  \item sin(t), cos(t): ℝ→[−1,1], Periode 2π
  \item cos(t) = sin(t+π/2)
  \item \textbf{Identität}: cos²(t)+sin²(t)=1
\end{itemize}


\subsubsection{Wellenparameter}

{\small
\begin{tabularx}{\linewidth}{XX}
\hline
\rowcolor{tableheadbg}
{\color{white}\textbf{Größe}} & {\color{white}\textbf{Formel}} \\[2pt]
\hline
\endhead
\hline
\endfoot
Wellenlänge λ & λ=1/f=2π/ω \\
Frequenz f & f=1/λ=ω/(2π) \\
Kreisfrequenz ω & ω=2πf=2π/λ \\
Amplitude A & maximale Auslenkung \\
\end{tabularx}
}


Standardform: A·sin(ωt) mit ω=2πf



\subsubsection{Ableitungs-Zyklus}
\begin{itemize}
  \item ∂/∂t [A·sin(ωt)] = Aω·cos(ωt)
  \item ∂/∂t [A·cos(ωt)] = −Aω·sin(ωt)
  \item sin → cos → −sin → −cos → sin
\end{itemize}


\subsection{Tangens / Sekans}
\begin{itemize}
  \item tan(t) = sin(t)/cos(t), Df = ℝ\textbackslash{}\{(1+2n)π/2\}
  \item ∂tan(t)/∂t = 1/cos²(t) = sec²(t) = 1+tan²(t)
\end{itemize}


\subsection{Kehrwertfunktionen}

{\small
\begin{tabularx}{\linewidth}{XX}
\hline
\rowcolor{tableheadbg}
{\color{white}\textbf{Name}} & {\color{white}\textbf{Definition}} \\[2pt]
\hline
\endhead
\hline
\endfoot
Kosekans csc(t) & 1/sin(t) \\
Sekans sec(t) & 1/cos(t) \\
Kotangens cot(t) & 1/tan(t) \\
\end{tabularx}
}


∂csc(t)/∂t = −cot(t)·csc(t)



\subsection{Umkehrfunktionen (arcus)}
\begin{itemize}
  \item arcsin, arccos, arctan — \textbf{nicht} sin⁻¹ (= Kehrwert!)
  \item Nicht eindeutig (Hauptzweig)
\end{itemize}


{\small
\begin{tabularx}{\linewidth}{XX}
\hline
\rowcolor{tableheadbg}
{\color{white}\textbf{Funktion}} & {\color{white}\textbf{Ableitung}} \\[2pt]
\hline
\endhead
\hline
\endfoot
arcsin(t) & 1/√(1−t²) \\
arccos(t) & −1/√(1−t²) \\
arctan(t) & 1/(t²+1) \\
\end{tabularx}
}



\subsection{Hyperbolische Trig.}
Hyperbel x²−y²=1 statt Einheitskreis.

\begin{itemize}
  \item sinh, cosh, tanh = sinh/cosh
  \item \textbf{Identität}: cosh²(t)−sinh²(t)=1
\end{itemize}


\subsection{Bewiesene Identitäten}
\begin{itemize}
  \item √(1−cos²(x)) = sin(x) für x∈(0,π)
  \item √(sec²(x)−1) = tan(x) für x∈(0, π/2)
\end{itemize}


\subsection{Beispielableitungen (trig.)}

{\small
\begin{tabularx}{\linewidth}{XXX}
\hline
\rowcolor{tableheadbg}
{\color{white}\textbf{f}} & {\color{white}\textbf{∂f/∂x}} & {\color{white}\textbf{∂f/∂y}} \\[2pt]
\hline
\endhead
\hline
\endfoot
sin(x)cos(y) & cos(x)cos(y) & −sin(x)sin(y) \\
√sin(x) + sin(√y) & cos(x)/(2√sin(x)) & cos(√y)/(2√y) \\
xy·e\textasciicircum{}sin(x) & y·e\textasciicircum{}sin(x)+xy·cos(x)·e\textasciicircum{}sin(x) & x·e\textasciicircum{}sin(x) \\
sin²(x)cos(x)+sin(ln(y)) & 2sin(x)cos²(x)−sin³(x) & cos(ln(y))/y \\
sin(arctan(x)) & 1/√(x²+1) & — \\
\end{tabularx}
}



\subsection{Data-Science-Bezug}
\begin{itemize}
  \item tanh(x): Aktivierungsfunktion in Neuronen
  \item cos(x): Kosinus-Distanz
  \item cos(ωt)+i·sin(ωt): Fourier-Analyse
\end{itemize}


\section{Komplexe Zahlen}


\subsection{Definition}
\textbf{Komplexe Zahlen ℂ}: Erweiterung von ℝ um die imaginäre Einheit i = √−1, um Wurzeln negativer Zahlen darzustellen.

\begin{itemize}
  \item ℂ = \{a+ib | a,b ∈ ℝ\}
  \item Re(z) = a, Im(z) = b
  \item Hierarchie: ℕ ⊂ ℕ₀ ⊂ ℤ ⊂ ℚ ⊂ ℝ ⊂ ℂ
\end{itemize}


\subsection{Darstellung}
\begin{itemize}
  \item \textbf{Kartesisch}: z = a+ib
  \item \textbf{Polar}: z = r·[cos(φ)+i·sin(φ)] = r·$e^{iφ}$
  \item \textbf{Betrag}: |z| = √(a²+b²); |z|² = Re(z)²+Im(z)²
  \item \textbf{Winkel}: φ = arctan(b/a)
\end{itemize}


\subsection{Konjugation}
z = a+ib → z̄ = a−ib



\subsection{Rechenoperationen}

{\small
\begin{tabularx}{\linewidth}{XX}
\hline
\rowcolor{tableheadbg}
{\color{white}\textbf{Operation}} & {\color{white}\textbf{Formel}} \\[2pt]
\hline
\endhead
\hline
\endfoot
Addition & z₁±z₂ = (a₁±a₂)+i(b₁±b₂) \\
Multiplikation & z₁z₂ = (a₁a₂−b₁b₂)+i(a₁b₂+a₂b₁) \\
Quadrierung & z² = (a²−b²)+2ab·i \\
Division & z₁/z₂ = (a₁a₂+b₁b₂)/(a₂²+b₂²)+i(a₂b₁−a₁b₂)/(a₂²+b₂²) \\
Skalar λ·z & λa+iλb; ·(−1) → Vorzeichenwechsel beider \\
\end{tabularx}
}


\begin{quote}
\textbf{Merke:} Multiplikation mit i = Drehung um 90°.
\end{quote}


\subsection{Eulersche Formel}
\textbf{$e^{iφ}$ = cos(φ)+i·sin(φ)} — verbindet Polarform mit Exponentialdarstellung.

\begin{itemize}
  \item z = a+ib = r·$e^{iφ}$
  \item A·[cos(ωt)+i·sin(ωt)] = A·$e^{iωt}$
  \item e\textasciicircum{}z = eˣ·[cos(y)+i·sin(y)]
  \item $e^{iπ}$+1 = 0
\end{itemize}


\subsection{Potenzieren}
\textbf{zˣ = (r·$e^{iφ}$)ˣ = rˣ·$e^{iφx}$} — erhöht Radius, dreht Winkel.



\subsection{Beispiele (a=−8+6i, b=12−9i, c=1+i)}
\begin{itemize}
  \item |a| = √(64+36) = 10
  \item |b|+Re(c) = 15+1 = 16
  \item ac−bc = c(a−b) = (1+i)(−20+15i) = −35−5i
  \item a+b+c² = 4−i (mit c² = 2i)
  \item a/b = −2/3
\end{itemize}


\section{Integralrechnung}


\subsection{Stammfunktion / Hauptsatz}
\textbf{Stammfunktion}: F mit F'(x)=f(x), F(x)=∫f(x)dx

\begin{itemize}
  \item ∫ₐᵇ f(x)dx = F(b)−F(a)
  \item Ober-/Untersumme: Näherung über Rechtecke (Max bzw. Min)
  \item Riemann-integrierbar: alle auf I stetigen Funktionen
\end{itemize}


\subsection{Standardintegrale}
\begin{itemize}
  \item ∫xᵇdx = xᵇ⁺¹/(b+1) für b≠−1
  \item ∫(1/x)dx = ln|x|
  \item ∫ln(x)dx = x·ln(x)−x
  \item ∫eˣdx = eˣ
  \item ∫f'(x)·e\textasciicircum{}f(x)dx = e\textasciicircum{}f(x)
\end{itemize}


\subsection{Partielle Integration}
\textbf{Formel}: ∫ₐᵇ f'(x)g(x)dx = [f(x)g(x)]ₐᵇ − ∫ₐᵇ f(x)g'(x)dx



\subsubsection{Beispiele}
\begin{itemize}
  \item ∫x·e⁻ˣdx = −(1+x)e⁻ˣ
  \item ∫x²eˣdx = x²eˣ−2xeˣ+2eˣ
  \item ∫sin(u)e⁻ᵘdu = −0.5·e⁻ˣ(sin(x)+cos(x)) (zweimal PI, Selbstreferenz)
  \item \textbf{Trick (Selbstreferenz)}: H = ∫ln(x)·x⁻¹dx → 2H = (ln x)² → H = 0.5(ln x)²
\end{itemize}


\subsection{Substitution (Umkehrung Kettenregel)}
\textbf{Formel}: ∫ₐᵇ f(g(x))·g'(x)dx = $\int_{g(a)}^{g(b)}$ f(u)du



\subsubsection{Schritte}
\begin{enumerate}
  \item u = g(x), du = g'(x)dx
  \item Einsetzen, integrieren
  \item Rücksubstitution u → g(x)
\end{enumerate}


\subsubsection{Beispiele}
\begin{itemize}
  \item ∫2x(x²+1)⁵dx, u=x²+1 → ⅙(x²+1)⁶
  \item ∫x³/(1+x⁴)dx, u=1+x⁴ → ½ln(1+x⁴) bzw. ½√(1+x⁴) je nach Form
  \item ∫x·ln(x²)dx → 0.5x²(ln(x²)−1)
  \item ∫−sin(x)·sin(cos(x))dx → −cos(cos(x))
  \item ∫1/(x·√(1+ln(x)))dx, u=1+ln(x) → 2√(1+ln(x))
  \item ∫(3x²−2x)·ln(x³−x²)dx, u=x³−x² → (x³−x²)(ln(x³−x²)−1)
  \item ∫0.5x·sin(x²)dx, u=x² → −0.25cos(x²)+C
\end{itemize}

\begin{quote}
\textbf{Merke}: ∫f(g(x))·g'(x)dx → u=g(x); konstante Faktoren vor das Integral.
\end{quote}


\subsection{Bestimmtes Integral — Beispiel}
\begin{itemize}
  \item ∫₁²(1/x²)dx = [−1/x]₁² = −0.5−(−1) = 0.5
\end{itemize}


\subsection{Trigonometrische Umformung}
∫[cos³(x)+sin³(x)/tan(x)]dx

\begin{enumerate}
  \item sin³(x)/tan(x) = sin²(x)·cos(x)
  \item cos(x)·[cos²+sin²] = cos(x)
  \item = sin(x)+C
\end{enumerate}


\subsection{Integrale über Skalarfelder (Doppelintegrale)}
\textbf{Doppelintegral}: ∫∫ f(x,y)dxdy — von innen nach außen integrieren.



\subsubsection{Strukturen}
\begin{itemize}
  \item Produkt: ∫∫fₓ(x)·fy(y)dxdy = (∫fy(y)dy)·(∫fₓ(x)dx)
  \item Additiv: ∫∫[fₓ(x)+fy(y)+c]dxdy → trennbar
\end{itemize}


\subsubsection{Beispiele}
\begin{itemize}
  \item ∫₀¹∫₀² xy dxdy = 1
  \item ∫₀¹∫₀² (xy+10) dxdy = 19
  \item ∫₀¹∫₀² (10−(x−1)²−y²) dxdy = 56/3
  \item ∫₀⁴∫₀¹ (x³y−5) dxdy = −18
  \item ∫₀⁴∫₀¹ (x²y−xy²) dxdy = −8
  \item ∫₀¹∫₁² (x+y+xy+1) dxdy = 3.75
  \item ∫∫(x+x/y) dxdy = x²·ln(y)+(2/3)y·x³+y·c₁+c₂
\end{itemize}


\subsection{Höhere Integralarten (nicht behandelt)}
\textbf{Satz von Stokes / Gauß}: ∫∫\_A div(F)d²A = ∮ F·n dS



\section{Gewöhnliche Differenzialgleichungen}


\subsection{Definition}
\textbf{DGL}: Gleichung, die Funktion mit mind. einer ihrer Ableitungen verknüpft. Lösung: Funktion y(x).



\subsection{Klassifikation}

{\small
\begin{tabularx}{\linewidth}{XXX}
\hline
\rowcolor{tableheadbg}
{\color{white}\textbf{Klasse}} & {\color{white}\textbf{Definition}} & {\color{white}\textbf{Beispiel}} \\[2pt]
\hline
\endhead
\hline
\endfoot
Gewöhnlich & nur eine Variable & y'(x)=… \\
Partiell & mehrere Variablen, partielle Abl. & ∂T/∂t=c·∂²T/∂x² \\
Ordnung n & höchste vorkommende Ableitung & y'(x)=… → 1. Ordnung \\
Explizit & nach höchster Ableitung aufgelöst & y'(x)=… \\
Linear & f und f' nur linear & y'=x²y \\
Nichtlinear & als Potenz/verschachtelt & y'=sin(y²) \\
Homogen & nur f und Ableitungen & y'=x²y \\
Inhomogen & zusätzliche Funktionen von x & y'=x²y+x \\
\end{tabularx}
}



\subsection{Linear homogen, 1. Ordnung}
\textbf{Form}: y'(x) = a(x)·y(x)

\textbf{Lösung allgemein}: y(x) = c₁·e\textasciicircum{}A(x), A(x)=∫a(x)dx

\textbf{Konstanter Koeffizient}: y'=αy → y(x) = c₁·$e^{αx}$

\textbf{AWP}: y(x) = y₀·$e^{α(x−x₀}$)



\subsubsection{Beispiele}

{\small
\begin{tabularx}{\linewidth}{XX}
\hline
\rowcolor{tableheadbg}
{\color{white}\textbf{AWP}} & {\color{white}\textbf{Lösung}} \\[2pt]
\hline
\endhead
\hline
\endfoot
y'=2y, y(4)=8 & y=8·$e^{2(x−4}$) \\
y'=ln(10)·y, y(0)=1 & y=10ˣ \\
y'=2x·y, y(0)=1 & y=$e^{x²}$ \\
y'=(1−x)y, y(2)=15 & y=15·$e^{x−0.5x²}$ \\
y'=y/x, y(1)=5 & y=5x \\
y'=y/2, y(1)=1 & y=$e^{0.5(x−1}$) \\
y'=ln2·y, y(0)=1 & y=2ˣ \\
y'=x²y, y(0)=4 & y=4·$e^{x³/3}$ \\
y'=eˣy, y(0)=2e & y=2e·$e^{eˣ}$ \\
\end{tabularx}
}



\subsubsection{Halbwertszeit / Zerfall}
\textbf{Zerfallsgesetz}: y'=−λy, λ=ln(2)/t₁/₂, y(t)=m₀·$e^{−λt}$; t₁/₂ = −ln(2)/α


\paragraph{Anwendung C14}\mbox{}\\[2pt]
y'(t)=αy(t), y(0)=1000, y(5700)=500 → α=ln(0.5)/5700=−0.0001216


\paragraph{Anwendung allgemein}\mbox{}\\[2pt]
y(0)=118.5, y(1)=23.7 → e\textasciicircum{}α=0.2 → α=ln(0.2)=−1.609 → t₁/₂=0.431 h=25.8 min



\subsection{Linear inhomogen, 1. Ordnung}
\textbf{Form}: y'(x) = a(x)·y(x)+b(x)

\textbf{Lösungsformel AWP}: y(x) = e\textasciicircum{}A(x)·(y₀ + $\int_{x₀}^{x}$ b(u)·$e^{−A(u}$)du), A(x)=∫a(x)dx



\subsubsection{Beispiele}

{\small
\begin{tabularx}{\linewidth}{XX}
\hline
\rowcolor{tableheadbg}
{\color{white}\textbf{AWP}} & {\color{white}\textbf{Lösung}} \\[2pt]
\hline
\endhead
\hline
\endfoot
y'=y/x+x², y(0)=0 & y=½x³ \\
y'=y/x+3x³, y(2)=0 & y=x⁴−8x \\
y'=y+eˣ, y(0)=0 & y=x·eˣ \\
y'=y/x+x, y(1)=0 & y=2x\textasciicircum{}1.5−2x \\
y'=y+x³eˣ, y(0)=0 & y=0.25·eˣ·x⁴ \\
y'=y/x−x, y(0)=0 & y=−2x²+x·c (c∈ℝ) \\
\end{tabularx}
}



\subsubsection{Anwendung CO₂-Raum}
200m³ Luft, +10g/min CO₂, 2m³/min Frischluft (1\%/min Abnahme)

\begin{itemize}
  \item DGL: y'=−0.01y+10, y(0)=0
  \item Lösung: y(x) = 1000−1000·$e^{−0.01x}$
\end{itemize}


\subsubsection{Anwendung Fliegen-Modell}
y'=−0.1y+10, y(0)=30 → y(x) = −70·$e^{−0.1x}$+100



\subsection{Nichtlinear: Getrennte Veränderliche}
\textbf{Form}: dy/dx = f(x)·g(y) (auch nichtlinear lösbar)



\subsubsection{Vorgehen}
\begin{enumerate}
  \item dy/g(y) = f(x)dx
  \item ∫dy/g(y) = ∫f(x)dx
  \item Nach y auflösen
  \item AWP einsetzen
\end{enumerate}


\subsubsection{Beispiele}

{\small
\begin{tabularx}{\linewidth}{XX}
\hline
\rowcolor{tableheadbg}
{\color{white}\textbf{AWP}} & {\color{white}\textbf{Lösung}} \\[2pt]
\hline
\endhead
\hline
\endfoot
y'=x·y², y(0)=1 & y=−2/(x²−2) \\
y'=x/y, y(0)=2 & y=√(x²+4) \\
y'=2x·e\textasciicircum{}y, y(0)=0 & y=−ln(1−x²) \\
y'=x·y³, y(0)=1 & y=1/√(1−x²) \\
y'=x·√y, y(0)=0 & y=x⁴/16 \\
y'=1/sin(y), y(0)=0 & y=arccos(1−x) \\
y'=y/tan(x), y(0)=0 & y=c₁·sin(x), c₁∈ℝ \\
y'=y+sin(x), y(0)=0 & y=−0.5(sin(x)+cos(x))+0.5eˣ \\
y'=−eʸ, y(0)=0 & y=−ln(x+1) \\
\end{tabularx}
}



\section{Fourieranalyse}


\subsection{Wellenüberlagerung (Linearität)}
\begin{itemize}
  \item Wellen beeinflussen sich nicht; Gesamtamplitude = Summe der Teilamplituden
  \item \textbf{Konstruktive Interferenz}: gleichphasig → Verstärkung
  \item \textbf{Destruktive Interferenz}: gegenphasig → Abschwächung
  \item Einzelwellen bleiben unverändert
\end{itemize}


\subsection{Fouriertransformation}
\textbf{Fouriertransformation}: Abbildung Zeit ↔ Spektrum.

\begin{itemize}
  \item f(t): Amplitude über Zeit
  \item |F(ω)|: Anteil der Frequenz am Gesamtsignal
\end{itemize}


{\small
\begin{tabularx}{\linewidth}{XX}
\hline
\rowcolor{tableheadbg}
{\color{white}\textbf{Richtung}} & {\color{white}\textbf{Formel}} \\[2pt]
\hline
\endhead
\hline
\endfoot
Hin & F(ω) = $\int_{−∞}^{∞}$ f(t)·$e^{−iωt}$ dt \\
Rück & f(t) = (1/2π)·$\int_{−∞}^{∞}$ F(ω)·$e^{iωt}$ dω \\
\end{tabularx}
}



\subsection{Diskrete Fouriertransformation (DFT)}
\begin{itemize}
  \item Für Zeitreihen
  \item Analytische Lösung von Hand nicht praktikabel
  \item Implementierung: \textbf{FFT} (Fast Fourier Transformation)
\end{itemize}


\subsection{Anwendung Data Science}
\begin{itemize}
  \item Periodische Komponenten in Zeitreihen identifizieren
\end{itemize}


\section{Klausurrahmen}


\subsection{Hilfsmittel}

{\small
\begin{tabularx}{\linewidth}{XX}
\hline
\rowcolor{tableheadbg}
{\color{white}\textbf{Hilfsmittel}} & {\color{white}\textbf{Zugelassen}} \\[2pt]
\hline
\endhead
\hline
\endfoot
Stifte, Lineale, Geodreiecke & ✓ \\
Skripte, Mathebücher, Formelsammlungen (gedruckt/handschriftlich) & ✓ \\
Taschenrechner & ✗ \\
Laptops, Tablets, E-Reader, Smartphones, Smartwatches & ✗ \\
\end{tabularx}
}



\subsection{Punkteverteilung}

{\small
\begin{tabularx}{\linewidth}{XXX}
\hline
\rowcolor{tableheadbg}
{\color{white}\textbf{Aufgabe}} & {\color{white}\textbf{Thema}} & {\color{white}\textbf{Punkte}} \\[2pt]
\hline
\endhead
\hline
\endfoot
1 & Vektoranalysis & 9 \\
2 & Gewöhnliche DGL & 8 \\
3 & Integralrechnung & 8 \\
\textbf{Gesamt} &  & \textbf{25} \\
\end{tabularx}
}





% ──────────────────────────────────────────────────────────────

\end{document}
